\bookchapter{第三十三章：未來展望——手機嘅下一程}

\chapterintro{手機唔會一夜之間消失，但佢嘅「獨霸地位」可能被新一代設備分流。呢一章我哋傾三件事：下一代設備生乜樣、手機點樣「隱身」為新身份基石、以及嗰啲我哋必須提前思考嘅倫理議題。}

\booksection{一、AR/VR與空間計算：手機之外嘅「第二桌面」}

下一代設備嘅競爭已經打響，最有可能分流手機時間嘅，係AR/VR與空間計算設備：

\begin{itemize}
  \item \textbf{Vision Pro後續}：蘋果Vision Pro嘅發布標誌住「空間計算」概念嘅正式落地，但第一代產品價格（29999蚊起）同重量仍係門檻。預計2026-2027年會推出更輕、更平嘅平價版Vision，等空間計算由「開發者玩具」走向「早期大眾」。屆時iPhone可能成為Vision眼鏡嘅算力伴侶與數據中樞。
  \item \textbf{AR眼鏡}：畀你眼前開個「第二桌面」，數字資訊變得好似空氣咁無處不喺。Meta與Ray-Ban合作嘅智能眼鏡已賣出數百萬副，國內Rokid、XREAL、INMO等廠商亦喺推出消費級AR眼鏡。2025-2026年嘅AR眼鏡已能做實時翻譯、導航HUD、AI助手對話，但仍受限於光波導顯示同續航，距離「取代手機」仲有3-5年。
  \item \textbf{智能戒指與可穿戴}：Samsung Galaxy Ring、Oura Ring等智能戒指開始流行，作為手機嘅「貼身傳感器」補充健康監測。手機嘅角色正由「萬能終端」轉向「中樞+多個輕量化配件」嘅組合。
\end{itemize}

\booksection{二、腦機介面：思想成為交互媒介}

思想本身將成為交互嘅媒介，但呢條路比想像中漫長：

\begin{itemize}
  \item \textbf{Neuralink最新進展}：馬斯克嘅Neuralink已於2024年完成首例人體植入，首位受試者Noland Arbaugh透過意念控制電腦滑鼠下國際象棋。目前植入式BCI（腦機介面）主要面向癱瘓患者輔助場景，目標係幫助佢哋重新「操作」數字世界。
  \item \textbf{非侵入式BCI}：多家公司開發非侵入式頭環（如Kernel、OpenBCI），無需手術即可讀取腦電信號，但精度遠低於侵入式。距離「用意念打字」仲有相當距離。
  \item \textbf{消費級展望}：腦機介面距離消費級手機交互仲有10-20年，但短期可能先用於無障礙輔助、專注力監測等場景。倫理與私隱問題（「我嘅腦電波算邊個嘅私隱」）將先於技術成熟成為爭論焦點。
\end{itemize}

\booksection{三、手機嘅「隱身」與新身份基石}

手機嘅另一個未來進化方向或者係「隱身」——化身為無處不喺嘅「幕後推手」，成為你數字身份嘅唯一基石：

\begin{itemize}
  \item \textbf{雲端雙因素認證}：瞬間完成「你係你」嘅認證，登入任何服務無需手動輸驗證碼。
  \item \textbf{數字遺產成為出廠設定}：新機激活時預設數字資產繼承方案，等「身後事」成為出廠預設。
  \item \textbf{由「設備中心」到「身份中心」}：你嘅雲端個人數據保險櫃先至係中心，硬件淨係訪問入口。換機、換設備，身份與數據始終跟住你。
\end{itemize}

\booksection{四、倫理思考：我哋，仲係手機嘅主人嗎？}

一個終極拷問始終存在：我哋係被技術賦能嘅「神」，仲係被算法餵養嘅「缸中之腦」？

\begin{itemize}
  \item \textbf{AI倫理}：當端側AI比你自己更懂你嘅偏好，你嘅「選擇」係自主意志定係被引導嘅結果？AI Agent代你落單、回覆消息時，責任與邊界喺邊度？
  \item \textbf{數據主權}：你嘅健康數據、腦電數據、生物特徵數據，係屬於你、屬於設備廠商仲係屬於雲服務商？《個人資訊保護法》同歐盟GDPR淨係起點，「數據可攜帶權」同「被遺忘權」嘅真正落地仲有好長嘅路。
  \item \textbf{數字鴻溝}：AI手機、摺疊屏、衞星通信都喺拉大「用得起」同「用唔起」之間嘅差距。數字基礎設施嘅公平接入，將成為新嘅社會議題——就好似20年前「識唔識上網」咁樣關鍵。
  \item 當算法比你更懂你想睇乜，你嘅「喜歡」有幾多係出於自發？
  \item 當你每隔幾分鐘就唔自覺咁點亮螢幕，呢個係你嘅習慣定係被設計出嚟嘅「上癮行為」？
\end{itemize}

\booksection{五、保持連接，亦保持獨立——你嘅數字生活哲學}

無論技術點樣進化，下面呢幾條數字生活哲學，10年後依然適用：

\begin{enumerate}
  \item 做資源嘅分配者，而唔係被推送嘅接收者。
  \item 等自動化成為你忠實嘅「副駕駛」。
  \item 建立你穩固嘅「數字後方」。
  \item 定期執行「數字斷捨離」與「社交齋戒」。
  \item 永遠將安全同私隱作為一切便利嘅前提。
\end{enumerate}

\begin{tcbpractice}{嚟試嚇：制定你嘅《個人數字設備使用服務條款》}
為你同你嘅手機（以及未來嘅AR眼鏡、AI助手）寫一份簡單嘅「服務條款」，嚟定義你哋嘅關係。例如：

\begin{itemize}
  \item 第1條：我嘅注意力時間屬於我。你喺我工作/陪伴家人時應主動進入「勿擾」狀態。
  \item 第2條：我嘅個人數據係我嘅財產。你嘅首要職責係安全加密與便捷同步。
  \item 第3條：我哋嘅共同目標係提升生活效率與體驗。若你透過誘導性設計等我沉迷，我有權將你卸載。
\end{itemize}

\end{tcbpractice}

\begin{tcbsummary}{本章小結}
\begin{itemize}
  \item Vision Pro與AR眼鏡係手機之外嘅新計算入口，但短期難以取代手機
  \item 腦機介面距離消費級仲有10-20年，倫理問題先於技術成熟
  \item 手機正朝「隱身」方向進化，由設備中心走向身份中心
  \item AI倫理、數據主權、數字鴻溝係未來10年嘅核心議題
  \item 保持連接亦保持獨立，係數字生活不變嘅哲學
\end{itemize}

\end{tcbsummary}
